\documentclass[11pt,a4paper]{article}

\usepackage[utf8]{inputenc}
\usepackage[T1]{fontenc}
\usepackage{textcomp}
\usepackage[french]{babel}
\addto\captionsfrench{\renewcommand{\tablename}{Tableau}
  \renewcommand{\figurename}{Figure}}

\usepackage{lmodern}\usepackage{geometry}\usepackage{setspace}
\usepackage{microtype}\usepackage{graphicx}\usepackage{xcolor}\usepackage{colortbl}
\usepackage{booktabs}\usepackage{tabularx}\usepackage{longtable}
\usepackage{array}\usepackage{enumitem}
\usepackage{titlesec}\usepackage{fancyhdr}\usepackage{caption}
\usepackage{float}\usepackage{needspace}
\usepackage{amssymb}\usepackage{amsmath}\usepackage[most]{tcolorbox}\usepackage{tikz}
\usetikzlibrary{positioning,arrows.meta,fit,shapes.geometric,calc}
\usepackage{adjustbox}
\usepackage{hyperref}
\geometry{top=2.1cm,bottom=2.1cm,left=2.2cm,right=2.1cm}
\setstretch{1.04}
\setlength{\emergencystretch}{1.2em}
\setlength{\parskip}{0.25em}

\definecolor{mainblue}{RGB}{28,74,102}
\definecolor{darkgray}{RGB}{62,62,66}
\definecolor{lightblue}{RGB}{238,245,250}
\definecolor{ent}{RGB}{221,236,247}
\definecolor{tab}{RGB}{215,237,221}
\definecolor{mod}{RGB}{233,233,238}
\definecolor{ext}{RGB}{250,233,207}
\definecolor{alerte}{RGB}{250,226,226}
\definecolor{codebg}{RGB}{246,246,248}
\definecolor{okgreen}{RGB}{46,125,50}
\definecolor{kored}{RGB}{176,42,42}

\hypersetup{colorlinks=true,linkcolor=mainblue,urlcolor=mainblue,
  pdftitle={Assistant bancaire frugal - resume de la proposition},
  pdfauthor={Equipe Yahya, Sohayb, Younes}}

\titleformat{\section}{\Large\bfseries\color{mainblue}}{\thesection}{0.6em}{}
\titleformat{\subsection}{\large\bfseries\color{darkgray}}{\thesubsection}{0.6em}{}
\titlespacing*{\section}{0pt}{1.2em}{0.45em}
\titlespacing*{\subsection}{0pt}{0.9em}{0.3em}

\captionsetup{font=small,labelfont=bf}
\setlist[itemize]{leftmargin=1.3em,itemsep=0.08em,topsep=0.2em}
\setlist[enumerate]{leftmargin=1.6em,itemsep=0.08em,topsep=0.2em}
\newcolumntype{Y}{>{\raggedright\arraybackslash}X}
\newcolumntype{L}[1]{>{\raggedright\arraybackslash}p{#1}}

\pagestyle{fancy}\fancyhf{}
\lhead{\small\color{darkgray}Assistant bancaire frugal --- résumé}
\rhead{\small\color{darkgray}Sujet 1}
\cfoot{\small\thepage}
\renewcommand{\headrulewidth}{0.3pt}

\newtcolorbox{cle}{breakable,colback=lightblue,colframe=mainblue,boxrule=0.7pt,arc=1.4mm,
  left=3mm,right=3mm,top=1.5mm,bottom=1.5mm}
\newtcolorbox{attention}{breakable,colback=alerte!60,colframe=kored!80,boxrule=0.6pt,arc=1.4mm,
  left=3mm,right=3mm,top=1.5mm,bottom=1.5mm}
\newtcolorbox{code}{breakable,colback=codebg,colframe=darkgray!35,boxrule=0.4pt,arc=1mm,
  left=2mm,right=2mm,top=0.8mm,bottom=0.8mm}

\DeclareRobustCommand{\f}[1]{{\texttt{\small
  \renewcommand{\_}{\char`\_\allowbreak}#1}}}
\newcommand{\eur}{\,\texteuro}
\newcommand{\amesurer}{\textit{\textcolor{darkgray}{à mesurer}}}

\tikzset{
  etage/.style={rectangle,draw=mainblue,fill=ent,rounded corners=2pt,
                align=center,font=\scriptsize,inner sep=4pt,minimum height=0.8cm},
  donnee/.style={rectangle,draw=mainblue!70,fill=tab,rounded corners=2pt,
                 align=center,font=\scriptsize,inner sep=4pt,minimum height=0.7cm},
  module/.style={rectangle,draw=darkgray!70,fill=mod,rounded corners=2pt,
                 align=center,font=\scriptsize,inner sep=4pt,minimum height=0.7cm},
  acteur/.style={rectangle,draw=darkgray,fill=ext,rounded corners=2pt,
                 align=center,font=\scriptsize,inner sep=4pt,minimum height=0.7cm},
  risque/.style={rectangle,draw=kored!80,fill=alerte,rounded corners=2pt,
                 align=center,font=\scriptsize,inner sep=4pt,minimum height=0.7cm},
  test/.style={diamond,aspect=2.2,draw=mainblue,fill=white,align=center,
               font=\scriptsize,inner sep=1pt},
  fl/.style={-{Latex[length=1.8mm]},semithick,draw=mainblue!80},
  flp/.style={-{Latex[length=1.8mm]},semithick,draw=darkgray!75,dashed},
  lien/.style={semithick,draw=darkgray!70},
  et/.style={font=\scriptsize,text=darkgray,inner sep=1.5pt},
}

\begin{document}

% ───────────────────────── En-tête ─────────────────────────
\begin{center}
{\small\scshape\color{darkgray} Hackathon IA --- La Banque Postale $\times$ CGI $\times$
ENSEIRB-MATMECA --- Sujet 1\par}
\vspace{0.25cm}
{\LARGE\bfseries\color{mainblue} Un assistant bancaire frugal, explicable, et qui agit\par}
\vspace{0.15cm}
{\large Résumé de la proposition de solution et de la répartition\par}
\vspace{0.2cm}
{\small Yahya (IA) --- Sohayb (génie logiciel) --- Younes (calcul intensif) --- lundi 26 octobre
2026\par}
{\footnotesize\color{darkgray} Version complète, avec annexes (85 intentions, catalogue, parcours,
glossaire, commandes) : \f{rapport/main.pdf}. Ce résumé fait foi là où il en diffère :
nous ne traduisons pas Banking77.\par}
\end{center}
\vspace{-0.2cm}

\section*{L'essentiel}

\textbf{Le problème.} Un client écrit à sa banque en langage libre (\og Je veux acheter un
appartement \fg, \og J'ai perdu ma carte \fg). L'assistant doit comprendre, poser les bonnes
questions, tenir compte de sa situation et proposer une piste --- avec une IA open source, locale,
la plus sobre possible, et dont chaque conseil se justifie.

\textbf{Notre réponse : une cascade.} Chaque étage ne passe la main au suivant que s'il ne sait pas
faire lui-même ; le calcul le plus cher n'arrive qu'en dernier, et c'est mesuré.

\begin{center}\small
\begin{tabularx}{\textwidth}{@{}l L{5.6cm} Y@{}}
\toprule
\textbf{Étage} & \textbf{Qui fait le travail} & \textbf{Pourquoi} \\
\midrule
1. Comprendre & Mots-clés, sinon \f{multilingual-e5-small} (118 M) + régression logistique &
  Moins d'une milliseconde pour les cas évidents ; mesurable sur le test officiel de Banking77 \\
2. Converser & Parcours YAML : quoi demander, dans quel ordre & On ne redemande jamais ce que le
  profil sait \\
3. Conseiller & Règles publiques (plafonds, 35\,\% d'endettement) + appétence & Une décision
  financière doit être reproductible et justifiable \\
4. Agir & Opposition, virement guidé, transmission au conseiller & Toujours après confirmation \\
5. Formuler & \textbf{Le seul appel LLM du tour, à la fin} : \f{ministral-3:3b} + garde-fou ;
  gabarits pour les réponses simples & Le LLM met en mots des faits calculés, n'en invente aucun \\
Escalade & \f{ministral-3:14b}, seulement si le doute persiste après une question & L'exception,
  comptée et affichée \\
\bottomrule
\end{tabularx}
\end{center}

\textbf{Quatre décisions.} (1) Comprendre sans LLM ; (2) décider sans LLM, avec les raisons
\emph{et} les produits écartés ; (3) un seul appel au plus petit LLM qui suffit, choisi en
descendant en taille depuis le 14B ; (4) tout mesurer contre une version \og tout LLM \fg.

\begin{cle}
\textbf{À retenir.} \emph{L'assistant comprend souvent avec de simples mots-clés, sinon avec 118
millions de paramètres ; il décide avec des règles que le client peut lire ; il n'appelle un LLM
qu'une fois, à la fin, pour mettre en mots ce qui est déjà calculé.} Un organisateur nous l'a
confirmé : \og pour gagner, c'est moins de LLM, plus d'algo \fg.
\end{cle}

% ───────────────────────── 1. Cadre ─────────────────────────
\section{Le cadre}

\paragraph{La règle commune.} Open source uniquement, aucune API de LLM externe, données en local ;
le plus petit modèle qui suffit, un plus gros seulement si le gain est \emph{mesuré} ; serveur GPU
sur site ; mesurer taille (en Go, quantisation comprise), énergie (CodeCarbon) et latence. Utiliser
d'autres IA \emph{pour développer} est autorisé. CGI : la frugalité, c'est \og surtout \fg{}
l'énergie ; un gros modèle \og à un instant précis \fg, c'est \og là que vous allez innover \fg.

\paragraph{Le sujet 1.} Objectifs : comprendre les demandes en langage naturel ; personnaliser
selon l'âge, les revenus et l'historique ; recommander un produit bancaire simple. Livrables :
\textbf{données préparées en français, modèle de compréhension des intentions, chatbot relié à une
base clients simulée, moteur de conseils personnalisés, mesures de frugalité}. Démarche attendue :
cadrer, préparer, concevoir, démontrer. À l'oral : on peut \emph{agir} (couper la carte) ou
seulement guider ; ils veulent notre regard de jeunes clients ; réponse immédiate.

\begin{center}\small
\begin{tabularx}{\textwidth}{@{}l L{5.2cm} L{3.1cm} Y@{}}
\toprule
\textbf{Thème} & \textbf{Message} & \textbf{Profil} & \textbf{Attendu} \\
\midrule
Carte & \og J'ai perdu ma carte hier soir, je fais quoi ? \fg & Inès, 34 ans & Réagir vite,
  opposition, rassurer \\
Frais & \og Pourquoi j'ai payé 8\eur{} de frais ce mois-ci ? \fg & Marc, 52 ans, découverts &
  Expliquer, orienter \\
Épargne & \og Je voudrais mettre de côté pour mes enfants. \fg & Sophie, 41 ans & Demander
  l'horizon \\
Virement & \og Je dois envoyer 500\eur{} à mon fils. \fg & Paul, 68 ans & Pas à pas, mots
  simples \\
Crédit & \og Je veux un crédit pour acheter une voiture. \fg & Karim, 29 ans, 1\,900\eur{} &
  Questionner, prudence \\
Insatisfaction & \og Votre appli ne marche jamais ! \fg & 3\textsuperscript{e} message du jour &
  Reconnaître, passer la main \\
Autres & Logement (énoncé), prêt étudiant et jargon (oral) & --- & Questionner, expliquer \\
\bottomrule
\end{tabularx}
\end{center}

\paragraph{Le jury.} Adrien (CGI), Denis Barthou (ENSEIRB), Jean-Michel (La Banque Postale, à
distance). 30 minutes : 15 de présentation \textbf{avec démo en direct sur le serveur fourni},
10 de questions, 5 de retour. Code déposé sur Git et support remis avant.

\begin{center}\small
\begin{tabular}{@{}l l l l l l@{}}
\toprule
Pertinence métier & Qualité technique & Frugalité & Explicabilité & Démonstration & Présentation \\
\midrule
\textbf{25\,\%} & \textbf{20\,\%} & \textbf{20\,\%} & \textbf{15\,\%} & \textbf{10\,\%} &
  \textbf{10\,\%} \\
\bottomrule
\end{tabular}
\end{center}

\paragraph{Ce que La Banque Postale valorise.} Sa priorité IA n\textsuperscript{o}\,1 est
l'expérience client (chatbot, réclamations, offres personnalisées). Elle insiste sur
l'accessibilité, l'inclusion des clients fragiles et la responsabilité. Un Français sur quatre est
son client.

\paragraph{Les conseils de l'organisateur.} Lors d'un échange lundi à 17h30 : moins de LLM, plus
d'algorithmes ou de modèles spécialisés ; en entrée, les mots-clés règlent environ 30\,\% des cas
pour presque rien et peuvent dire quand ils ne sont pas sûrs ; en sortie, un LLM, appelé une seule
fois, à la fin ; choisir sa taille en descendant (15B ou 30B, puis 4B, 2B, 1B) ; un seul sujet ;
d'abord une solution sûre, ensuite les idées exotiques ; se faire plaisir.

\paragraph{Les moyens.} Postes de la salle : RTX 3060 12 Go, Ollama avec \f{gemma4:e2b} (7,2 Go)
et \f{ministral-3:14b} (9,1 Go). DeepEirb : 2 $\times$ A6000 48 Go, SSH et vLLM \emph{sur demande
à Denis Barthou}. Jetson et OrangePi dès mardi, pour le déploiement. Un jour et demi de travail.

% ───────────────────────── 2. Données ─────────────────────────
\section{Les données}

\paragraph{Banking77.} 13\,083 questions de clients d'une néobanque britannique, \textbf{en
anglais}, 77 intentions ; 10\,003 en entraînement (35 à 187 par intention), 3\,080 en test (40 par
intention). Références de l'article : 87,19\,\% (BERT figé) à 93,66\,\% (BERT affiné) ; un
encodeur figé et un classifieur léger se placent à un ou deux points. \textbf{Il couvre le service
(carte, frais, virements), pas le conseil} : aucune demande d'épargne, de crédit ou de logement.
Il contient aussi un peu de bruit d'étiquetage. Nous rangeons les 77 intentions en 9 familles et
créons 8 intentions françaises :

\begin{center}\small
\begin{tabularx}{\textwidth}{@{}L{3.4cm} Y@{}}
\toprule
\textbf{Familles Banking77} & carte perdue, usage de la carte, commande de carte, frais et
  contestations, virements, compte et identité, espèces, change, alimentation du compte \\
\textbf{Nouvelles intentions} & \f{epargne\_projet}, \f{credit\_immobilier}, \f{credit\_conso},
  \f{pret\_etudiant}, \f{faire\_virement}, \f{question\_definition}, \f{insatisfaction},
  \f{hors\_sujet} \\
\textbf{Au total} & 14 familles, 85 intentions ; environ 80 exemples générés hors ligne par
  intention nouvelle (relus), 20 exemples de test écrits à la main \\
\bottomrule
\end{tabularx}
\end{center}

\paragraph{Pas de traduction de Banking77.} C'est notre choix. L'encodeur
\f{multilingual-e5-small} place une phrase et sa traduction très près l'une de l'autre dans
l'espace des vecteurs : un classifieur entraîné sur les phrases anglaises reconnaît les mêmes
demandes écrites en français. Traduire 13\,083 phrases coûterait du temps et de l'énergie, et
ajouterait des erreurs de traduction, pour un gain non prouvé. \textbf{Ce qui est en français} :
le lexique de mots-clés ; les exemples des 8 nouvelles intentions ; 5 phrases par famille
Banking77 ajoutées à l'entraînement ; un \textbf{test français de 250 phrases écrites à la main}
(10 par famille Banking77, 20 par nouvelle intention) ; le glossaire, le catalogue et les libellés
des profils. C'est notre livrable \og données préparées en français \fg. L'écart de précision
entre le test anglais officiel et notre test français mesure ce que l'encodeur transfère d'une
langue à l'autre : un résultat à présenter.

\paragraph{Bank Marketing.} 45\,211 clients d'une banque portugaise (2008--2010) : âge, métier,
situation, études, solde (3\,766 négatifs, 8,3\,\%), crédit immobilier (55,6\,\%), crédit personnel
(16,0\,\%) ; cible \f{y} = a souscrit un dépôt à terme (11,7\,\%). Pas de revenus, d'enfants ni
d'historique : on les simule (graine fixe, drapeau \f{simule=True}, affichés en italique) derrière un
adaptateur unique, prêt pour la version enrichie promise par les organisateurs.

\begin{attention}
\textbf{\f{duration} est une fuite} : la durée de l'appel n'est connue qu'après l'appel et
détermine presque la réponse. On l'exclut (avec \f{contact}, \f{day}, \f{month}). Le dire au jury.
\end{attention}

\paragraph{Les 8 clients de démonstration.} Ancrés sur des lignes réelles de Bank Marketing :
C01 Inès (carte), C02 Marc (frais, découverts), C03 Sophie (épargne, 2 enfants), C04 Paul (68 ans,
virement), C05 Karim (1\,900\eur{}, crédit en cours), C06 Julie (3 contacts aujourd'hui),
C07 Léa (2\,900\eur{}, 25\,000\eur{} d'épargne, appartement), C08 Lucas (20 ans, étudiant).

\paragraph{Le catalogue.} En JSON, une fiche par produit, avec \textbf{source publique et date} :
Livret A (22\,950\eur{}), LDDS (12\,000\eur{}), LEP (10\,000\eur{}, sous condition de revenu),
Livret Jeune (12--25 ans, 1\,600\eur{}), PEL (61\,200\eur{}), CEL (15\,300\eur{}), crédit
immobilier, PTZ, prêt personnel, prêt étudiant, offre clientèle fragile, opposition de carte.
\textbf{Aucun taux en dur} : c'est un paramètre daté du catalogue.

% ───────────────────────── 3. Solution ─────────────────────────
\section{La solution}

\begin{figure}[tbp]
\centering
\begin{adjustbox}{max width=0.92\textwidth}
\begin{tikzpicture}[x=1cm,y=1cm]
  \node[acteur,minimum width=3.6cm] (client) at (0,0) {Client, dans l'interface};
  \node[acteur,minimum width=5.4cm,text width=5.2cm] (out) at (0,-9.6)
    {Réponse\\+ panneau \og pourquoi \fg{} + panneau \og frugalité \fg};
  \node[etage,minimum width=4.4cm] (e1) at (0,-1.5)
    {\textbf{1. Comprendre}\\ mots-clés, sinon \f{multilingual-e5-small}\\ + régression logistique ; entités par règles};
  \node[test] (d) at (0,-3.05) {confiance $\geq \tau$ ?};
  \node[etage,minimum width=4.4cm] (e2) at (0,-4.65)
    {\textbf{2. Converser}\\ parcours YAML, profil du client};
  \node[etage,minimum width=3.3cm] (e3) at (-2.55,-6.35)
    {\textbf{3. Conseiller}\\ règles publiques\\ + appétence};
  \node[etage,minimum width=3.3cm] (e4) at (2.55,-6.35)
    {\textbf{4. Agir}\\ outils\\ + confirmation};
  \node[etage,minimum width=4.4cm] (e5) at (0,-8.05)
    {\textbf{5. Formuler}\\ gabarit, sinon \f{ministral-3:3b}\\ + garde-fou};
  \node[module,minimum width=2.9cm] (clar) at (-5.4,-3.05) {Question de\\ clarification};
  \node[risque,minimum width=2.9cm] (esc) at (-5.4,-4.65) {\textbf{Escalade}\\ \f{ministral-3:14b}};
  \node[donnee,minimum width=2.6cm] (prof) at (-5.6,-6.35) {Profils clients\\ (simulés)};
  \node[donnee,minimum width=2.6cm] (cata) at (-5.6,-7.6) {Catalogue\\ + sources};
  \node[donnee,minimum width=2.6cm] (jour) at (5.6,-6.35) {Journal\\ des actions};
  \node[donnee,minimum width=2.6cm] (glos) at (5.6,-8.05) {Glossaire\\ (fiches vérifiées)};
  \node[module,dashed,minimum width=3.2cm,text width=3.1cm] (off) at (5.95,-1.5)
    {\textit{Hors ligne, une fois}\\ exemples français des\\ 8 nouvelles intentions,\\ qui entraînent l'étage 1\\ (gros modèle, DeepEirb)};
  \node[module,minimum height=9.6cm,minimum width=0.75cm] (mes) at (8.15,-4.8) {};
  \node[rotate=90,font=\scriptsize,align=center] at (mes.center)
    {\textbf{Mesure continue} : latence, tokens, énergie, Go chargés};
  \draw[fl] (client) -- (e1);
  \draw[fl] (e1) -- (d);
  \draw[fl] (d) -- node[et,right]{oui} (e2);
  \draw[fl] (d) -- node[et,above]{non} (clar);
  \draw[flp] (clar.north) |- node[et,pos=0.25,left]{question} (client.west);
  \draw[flp] (clar) -- node[et,right]{si toujours ambigu} (esc);
  \draw[fl] (esc) -- (e2);
  \draw[fl] (e2) -- (e3);
  \draw[fl] (e2) -- (e4);
  \draw[fl] (e2) -- node[et,fill=white]{service simple} (e5);
  \draw[fl] (e3) -- (e5);
  \draw[fl] (e4) -- (e5);
  \draw[fl] (e5) -- (out);
  \draw[lien] (prof) -- (e3);
  \draw[lien] (cata) -- (e3.south west);
  \draw[lien] (e4) -- (jour);
  \draw[lien] (e5) -- (glos);
  \draw[flp] (off.west) -- (e1.east);
\end{tikzpicture}
\end{adjustbox}
\caption{La cascade. En bleu, sans gros modèle ; en rouge, l'exception comptée.}
\end{figure}

\subsection{Étage 1 --- Comprendre (cible : moins de 30 ms)}
\begin{itemize}
  \item \textbf{Mots-clés d'abord.} Lexique d'expressions par famille (\og perdu ma carte \fg,
        \og virement \fg, \og appartement \fg\dots), cherché dans le message normalisé (minuscules,
        sans accents). Il ne conclut que s'il est sûr : une seule famille touchée, phrase courte ou
        plusieurs expressions. Expressions gardées seulement si leur précision dépasse 97\,\%.
  \item \textbf{Sinon, classifieur à deux niveaux} (14 familles, puis 85 intentions) :
        \f{multilingual-e5-small} (préfixe \f{query:}) + régression logistique scikit-learn,
        entraînée sur Banking77 en anglais, \textbf{sans traduction}, plus nos phrases françaises. Quelques secondes
        d'entraînement sur CPU ; affiche aussi la phrase d'entraînement la plus proche.
  \item \textbf{Seuil $\tau$} choisi par balayage (97\,\% de précision sur les messages acceptés).
        En dessous : question fermée (\og votre carte, ou un virement ? \fg), puis escalade.
  \item \textbf{Contexte} : si un parcours attend un montant, \og 250\,000 \fg{} est une réponse.
        \textbf{Entités} par expressions régulières (montants, durées, liens, objets, dates) ; ce qui
        manque est demandé au client, pas deviné par un LLM. \textbf{Insatisfaction} :
        classifieur + nombre de contacts du jour.
\end{itemize}

\subsection{Étage 2 --- Converser}
Un fichier YAML par parcours : informations nécessaires, source (profil ou question), ordre, fin
(conseil, action, passage de main). Un automate déterministe les lit. \textbf{On confirme ce que
le profil sait au lieu de le demander.} Mode \og mots simples \fg{} automatique dès 70 ans ou
\og peu à l'aise en ligne \fg{} (phrases courtes, une étape par message). Glossaire : fiches
écrites et vérifiées par nous, retrouvées par similarité en mémoire (numpy), sans LLM.

\subsection{Étage 3 --- Conseiller}
Éligibilité (âge, revenu, plafonds), puis pertinence (horizon, objectif), puis classement. Chaque
règle a un identifiant et une phrase d'explication : la sortie donne \textbf{les propositions avec
leurs raisons, et les produits écartés avec la règle qui les écarte}. Appétence épargne : régression
logistique sur Bank Marketing, sans \f{duration}, évaluée en AUC, trois facteurs principaux
affichés ; LightGBM seulement si le gain est mesuré. Crédit immobilier (règles HCSF) :
$m = 0{,}35 \times R - C$ et capacité $= m \times \frac{1-(1+i)^{-n}}{i}$, durée $\leq$ 25 ans,
taux indicatif daté. Crédit : mention légale obligatoire, aucune promesse, rendez-vous conseiller.

\subsection{Étage 4 --- Agir}
\f{consulter\_profil} et \f{expliquer\_frais} (lecture libre) ; \f{bloquer\_carte},
\f{preparer\_virement} et \f{transmettre\_conseiller} (\textbf{après confirmation}, avec trace dans
un journal). Le conseiller reçoit un résumé des faits du parcours.

\subsection{Étage 5 --- Formuler}
\textbf{Un seul appel LLM par tour, à la fin}, après que tout a été compris, décidé et calculé.
Gabarits relus pour les réponses simples. Le petit LLM ne reçoit que la consigne, les faits
structurés et la question en cours. \textbf{Garde-fou} : chaque nombre du texte doit figurer dans
les faits, aucune formule interdite (\og accordé \fg, \og garanti \fg), longueur bornée ; sinon, le
gabarit s'affiche et l'échec est compté.

\subsection{Un exemple : l'appartement de Léa}
\begin{center}\small
\begin{tabularx}{\textwidth}{@{}l Y l@{}}
\toprule
\textbf{Tour} & \textbf{Ce qui se passe} & \textbf{Par} \\
\midrule
\og Je veux acheter un appartement \fg & Mots-clés \og acheter \fg{} + \og appartement \fg{} : une
  seule famille, phrase courte ; revenus connus, donc confirmés ; on demande le prix et l'apport &
  mots-clés, gabarit \\
\og Autour de 250\,000, je mets les 25\,000 \fg & Notaire 7,5\,\% $\Rightarrow$ besoin
  243\,750\eur{} ; à 3,5\,\% sur 25 ans, mensualité $\approx$ 1\,220\eur{}, soit
  \textbf{42,1\,\% $>$ 35\,\%} ; prix maximal $\approx$ 212\,000\eur{} ; leviers et rendez-vous ;
  écarté : prêt sur 30 ans & règles, \f{ministral-3:3b} \\
\bottomrule
\end{tabularx}
\end{center}

\paragraph{Conformité.} Données simulées et traitement local. Le règlement européen sur l'IA classe
à haut risque la notation de solvabilité : nous n'en faisons pas, un conseiller décide ;
l'utilisateur est informé qu'il parle à une IA. Rappel anti-arnaque : aucun code ne se communique.

% ───────────────────────── 4. Modèles et technologies ─────────────────────────
\section{Modèles et technologies}

\begin{center}\small
\begin{tabularx}{\textwidth}{@{}L{3.2cm} L{4.6cm} l Y@{}}
\toprule
\textbf{Rôle} & \textbf{Modèle} & \textbf{Taille} & \textbf{Où} \\
\midrule
Filtre d'entrée & Lexique de mots-clés & quelques Ko & CPU \\
Encodeur & \f{multilingual-e5-small} (118 M) & $\approx$\,0,47 Go & CPU ou GPU \\
Têtes famille + intention & Régression logistique ($\approx$\,38\,000 poids) & $<$ 1 Mo & CPU \\
Formulation & \f{ministral-3:3b} Q4\_K\_M (hypothèse) & 3,0 Go & RTX 3060 \\
Escalade & \f{ministral-3:14b} Q4\_K\_M & 9,1 Go & DeepEirb (vLLM) ou 8B en local \\
Appétence & Régression logistique & $<$ 1 Mo & CPU \\
Exemples des 8 intentions & \f{ministral-3:14b} via vLLM, relus à la main & --- & DeepEirb, une
  fois \\
\bottomrule
\end{tabularx}
\end{center}

\begin{cle}
\textbf{Choisir le LLM en descendant.} Référence \f{ministral-3:14b} (ou un 30B sur DeepEirb), puis
\f{gemma4:e2b} (7,2 Go), \f{ministral-3:8b} (6,0), \f{ministral-3:3b} (3,0), \f{qwen3:1.7b} (1,4,
réflexion coupée), \f{gemma3:1b} (0,8), \f{qwen3:0.6b} (0,5), \f{gemma3:270m} (0,3). Seuil fixé
avant : zéro nombre inventé et 27 réponses sur 30 jugées bonnes. \textbf{On garde le plus petit
qui tient le seuil} ; la courbe qualité--taille--énergie va dans le support.
\end{cle}

\begin{center}\small
\begin{tabularx}{\textwidth}{@{}L{3.6cm} Y@{}}
\toprule
\textbf{Brique} & \textbf{Technologie} \\
\midrule
Langage, environnement & Python 3, \f{venv}, Git (branches courtes) \\
Compréhension & \f{sentence-transformers}, \f{scikit-learn}, numpy, expressions régulières \\
Dialogue & PyYAML (parcours), automate maison \\
LLM & Ollama (API \f{/api/chat}, \f{num\_ctx} 4\,096, température 0,2, 200 tokens max,
  \f{keep\_alive}), vLLM sur DeepEirb (API compatible OpenAI) \\
Données & pandas (DuckDB si besoin), JSON pour catalogue et glossaire, CSV pour les clients \\
Interface & Streamlit : conversation, profil, panneaux \og pourquoi \fg{} et \og frugalité \fg,
  mode mots simples, taille du texte \\
Mesure & CodeCarbon (\f{OfflineEmissionsTracker}, pays FRA), NVML / \f{nvidia-smi},
  \f{tegrastats} sur Jetson \\
Tests & pytest pour les règles ; rejeu YAML des 30 scénarios \\
\bottomrule
\end{tabularx}
\end{center}

\textbf{Mémoire.} Le 3B et le 14B ne tiennent pas ensemble dans 12 Go : escalade sur DeepEirb, ou
8B en local, ou un seul modèle chargé à la fois.

% ───────────────────────── 5. Évaluation ─────────────────────────
\section{Évaluer et mesurer}

\begin{itemize}
  \item \textbf{Compréhension}, sur deux jeux : le test officiel Banking77 en anglais (3\,080,
        comparable à l'article) et notre test français écrit à la main (250) ; leur écart mesure le
        transfert de l'anglais au français. Précision et F1 macro, par intention et par famille ;
        part acceptée par les mots-clés (l'organisateur l'estime à 30\,\%). Comparé : mots-clés,
        TF-IDF, MiniLM, e5-small, e5-base, et deux LLM qui classent seuls (1.7B, 14B).
  \item \textbf{Scénarios} : 30 conversations YAML (9 parcours $\times$ 3 variantes + 3 hors
        périmètre). Vérifiés : famille, questions posées \emph{et interdites}, produits proposés et
        écartés, action, passage de main, aucun nombre inventé. Un scénario qui casse bloque la
        fusion.
  \item \textbf{Frugalité} : répartition par niveau (mots-clés, classifieur, escalade), part sans
        LLM, latence p50/p95, Go chargés, énergie par requête.
\end{itemize}

\begin{center}\small
\begin{tabular}{@{}l l c c c c c@{}}
\toprule
& \textbf{Configuration} & \textbf{Réussite} & \textbf{Wh / 1\,000 req.} & \textbf{Latence p95} &
  \textbf{Go} & \textbf{Sans LLM} \\
\midrule
A & Tout LLM : \f{ministral-3:14b} fait tout & \amesurer & \amesurer & \amesurer & 9,1 & 0\,\% \\
B & \textbf{Notre cascade} & \amesurer & \amesurer & \amesurer & 3,5 & \amesurer \\
C & Cascade, gabarits seuls & \amesurer & \amesurer & \amesurer & 0,5 & 100\,\% \\
\bottomrule
\end{tabular}
\end{center}

\textbf{Énergie} : CodeCarbon hors ligne, par lots de 200 requêtes, après chauffe, puissance de
repos (60 s) soustraite, trois répétitions : $(E - P_0 t)/N$ par requête. Le coût de la préparation
hors ligne est donné à part.

% ───────────────────────── 6. Démo ─────────────────────────
\section{La démonstration (6 minutes)}

\begin{center}\small
\begin{tabularx}{\textwidth}{@{}l L{3.2cm} Y l@{}}
\toprule
\textbf{Temps} & \textbf{Scène} & \textbf{Ce que le jury voit} & \textbf{Critère} \\
\midrule
0 min & Inès, carte perdue & Confirmation, opposition, référence ; compris par mots-clés, aucun
  LLM & Métier, frugalité \\
1 min & Léa, appartement & Revenus non redemandés ; refus chiffré ; pourquoi et pourquoi pas &
  Explicabilité \\
2 min 30 & Sophie puis Lucas & Même phrase, deux clients, deux conseils & Personnalisation \\
3 min 30 & Marc, 8\eur{} de frais & Historique, plafond légal, offre clientèle fragile & Inclusion \\
4 min 30 & Julie, mécontente & Excuses, passage de main avec résumé & Métier \\
5 min 15 & Bilan & Panneau frugalité cumulé ; graphique A/B/C & Frugalité \\
\bottomrule
\end{tabularx}
\end{center}

Secours : vidéo enregistrée mercredi matin, mode gabarits, vérification réelle d'Ollama, modèles
préchargés, messages de démo cliquables, répétition sur la machine de démo.

% ───────────────────────── 7. Organisation ─────────────────────────
\section{L'organisation}

\paragraph{Principes.} Un module, un responsable, un suppléant. Contrats figés lundi soir, bouchons
pour avancer sans attendre. \f{main} tourne toujours. Deux intégrations mardi (12h30, 18h), jugées
par les scénarios. D'abord une solution sûre (jusqu'à J2), ensuite les idées exotiques.

\begin{center}\small
\begin{tabularx}{\textwidth}{@{}L{2.4cm} Y Y Y@{}}
\toprule
& \textbf{M1 --- Comprendre et mesurer} & \textbf{M2 --- Converser et formuler} &
  \textbf{M3 --- Conseiller et agir} \\
\midrule
\textbf{Responsable} & \textbf{Younes} (suppléant Yahya) & \textbf{Yahya} (suppléant Sohayb) &
  \textbf{Sohayb} (suppléant Younes) \\
\textbf{Données} & Test français écrit à la main ; 8 intentions ; fiche de données & Glossaire, gabarits &
  8 clients simulés ; catalogue \\
\textbf{Modèle} & Mots-clés ; encodeur et têtes ; seuil $\tau$ & Choix du LLM en descendant ;
  garde-fou ; escalade & Appétence (AUC) \\
\textbf{Logiciel} & \f{comprendre()} ; instrumentation ; bancs NLU et A/B/C ; Jetson (bonus) &
  \f{reply()} ; automate et 9 parcours ; interface Streamlit ; vidéo de secours & Moteur de règles
  et calculs ; 5 outils ; rejeu des scénarios ; dépôt Git final \\
\textbf{Fichiers} & \f{src/comprendre/}, \f{src/mesure/}, \f{eval/bench\_*.py} &
  \f{src/assistant.py}, \f{src/converser/}, \f{app.py} & \f{src/conseiller/},
  \f{eval/run\_scenarios.py}, \f{tests/} \\
\textbf{Dans le pitch} & Résultats et frugalité & Problème, architecture, démo &
  Explicabilité et conformité \\
\bottomrule
\end{tabularx}
\end{center}

Chacun écrit au moins dix scénarios et la diapositive de son module. Le découpage suit la cascade,
pas les filières : chaque module a une part données, une part modèle et une part logiciel.

\paragraph{Le contrat commun.} Sohayb l'a posé dans \f{src/contrats.py} (commit du lundi 18h23),
avec le squelette de \f{reply()} et de l'interface. À relire ensemble ce soir, puis à figer :

\begin{code}
\footnotesize
\begin{verbatim}
Comprehension(famille, intention, confiance, alternatives, entites, insatisfaction,
              modele, exemple_proche)                                            # M1
Conseil(propositions=[Proposition(produit_id, raisons, sources, score)],
        ecartees, manque, alertes, faits)                                         # M3
Action(outil, parametres, confirmee, resultat)                                    # M3
Etape(nom, modele, latence_ms, tokens)                                            # M1
Reponse(texte, comprehension, conseil, action, escalade, trace=[Etape])           # M2
EtatConversation(intention, slots, derniere_reponse, historique, action_en_attente)
reply(message, client_id, etat) -> (Reponse, EtatConversation)   # point d'entree, M2
\end{verbatim}
\end{code}

\paragraph{Le dépôt.} \f{app.py} ; \f{src/} (\f{contrats.py}, \f{assistant.py},
\f{comprendre/}, \f{converser/}, \f{conseiller/}, \f{mesure/}) ; \f{data/raw/} (données téléchargées, versionnées) et
\f{data/processed/} (versionné : c'est un livrable) ; \f{eval/} (scénarios, bancs) ; \f{tests/}.

\begin{table}[H]
\centering\scriptsize
\renewcommand{\arraystretch}{1.2}
\begin{tabularx}{\textwidth}{@{}L{1.3cm} Y Y Y Y Y@{}}
\toprule
& \textbf{Lundi 18h--20h} & \textbf{Mardi 8h30--12h30} & \textbf{Mardi 13h30--18h}
  & \textbf{Mardi soir (option)} & \textbf{Mercredi 8h30--12h} \\
\midrule
\textbf{Younes} & \cellcolor{ent}Accès DeepEirb ; lancer la génération des 8 intentions
  & \cellcolor{ent}Mots-clés ; têtes v1 ; entités ; \texttt{comprendre()}
  & \cellcolor{ent}Test français ; 8 intentions ; $\tau$ ; escalade ; instrumentation
  & \cellcolor{ent}Banc A/B/C
  & \cellcolor{ent}Chiffres finaux ; graphique ; Jetson \\
\textbf{Yahya} & \cellcolor{tab}Parcours carte sur le squelette
  & \cellcolor{tab}Automate ; carte et frais ; gabarits ; LLM branché
  & \cellcolor{tab}Épargne, crédits, virement, insatisfaction ; garde-fou ; panneaux
  & \cellcolor{tab}Scénarios en échec
  & \cellcolor{tab}Vidéo de secours ; finitions \\
\textbf{Sohayb} & \cellcolor{ext}8 clients ; catalogue v0
  & \cellcolor{ext}\texttt{clients.py} ; règles épargne et crédit ; opposition ; frais
  & \cellcolor{ext}Appétence ; écartés ; virement ; conseiller ; scénarios
  & \cellcolor{ext}Scénarios en échec
  & \cellcolor{ext}Dépôt Git final ; README \\
\textbf{Tous} & \cellcolor{mod}Relire et figer le socle de Sohayb & \cellcolor{mod}Test français :
  environ 85 phrases chacun & & & \cellcolor{mod}Diapositives ; 2 répétitions \\
\midrule
\textbf{Jalon} & \textcolor{kored}{\textbf{J0}} 20h : \texttt{app.py} répond &
  \textcolor{kored}{\textbf{J1}} 12h30 : Inès et Marc de bout en bout &
  \textcolor{kored}{\textbf{J2}} 18h : 9 parcours, \textbf{gel} & --- &
  \textcolor{kored}{\textbf{J3}} 11h30 : code, vidéo, support remis \\
\bottomrule
\end{tabularx}
\caption{Le planning.}
\end{table}

\paragraph{Priorités.} \textbf{Indispensable à J2} : étage 1 mesuré ; parcours carte, frais,
épargne, appartement, crédit auto ; règles avec raisons et écartés ; 8 clients ; LLM et garde-fou ;
panneaux ; banc A/B/C. \textbf{Important} : virement en mots simples, insatisfaction, jargon, prêt
étudiant, appétence, escalade. \textbf{Si le temps le permet} : idées exotiques, Jetson, langue du
client, MCP, voix locale. \textbf{Hors périmètre} : affinage de LLM, vrai SI bancaire.

\paragraph{Règles.} Branches \f{m1/}, \f{m2/}, \f{m3/} fusionnées dans la demi-journée après relecture
du suppléant ; ni données brutes ni secrets dans Git ; point de cinq minutes toutes les trois heures ;
chaque chiffre noté avec sa commande. \textbf{\og Fini \fg} = dans \f{main}, scénarios verts,
mesures notées, diapositive écrite.

% ───────────────────────── 8. Risques ─────────────────────────
\section{Les risques principaux}

\begin{center}\small
\begin{tabularx}{\textwidth}{@{}L{4.6cm} Y l@{}}
\toprule
\textbf{Risque} & \textbf{Parade} & \textbf{Qui} \\
\midrule
DeepEirb inaccessible lundi soir & Exemples des 8 intentions générés par \f{ministral-3:14b} sur un
  poste de la salle ; escalade en 8B local & M1 \\
Transfert anglais $\to$ français insuffisant & Plus de phrases françaises écrites à la main dans
  l'entraînement ; les mots-clés en premier ; écart mesuré et présenté & M1 \\
Petit LLM faible en français & Gabarits ; la descente s'arrête plus haut (8B) si mesuré & M2 \\
Mots-clés trop sûrs d'eux & Précision $\geq$ 97\,\% par expression ; une seule famille ; sinon
  classifieur & M1 \\
Base enrichie en retard & Simulation dès lundi ; un seul adaptateur à changer & M3 \\
Intégration tardive & Contrats J0, bouchons, deux intégrations datées & Tous \\
Démo qui plante & Vidéo, mode gabarits, vérification d'Ollama, répétition sur place & M2 \\
\textbf{Le temps} & \textbf{Priorités figées ; gel à J2 ; on coupe sans débat} & Tous \\
\bottomrule
\end{tabularx}
\end{center}

% ───────────────────────── 9. Pitch ─────────────────────────
\section{Le pitch (15 minutes)}

\begin{center}\small
\begin{tabularx}{\textwidth}{@{}l L{3.2cm} Y L{2.4cm}@{}}
\toprule
\textbf{Temps} & \textbf{Partie} & \textbf{Contenu} & \textbf{Qui} \\
\midrule
0 min & Le problème & Client, priorité de La Banque Postale, règle commune & Yahya \\
1 min 30 & La démarche & Quatre temps, constat Banking77, la cascade & Yahya \\
4 min & La démo & Le script ci-dessus & Yahya pilote, Sohayb commente \\
10 min & Les résultats & Tests anglais et français, scénarios, A/B/C, coût de la préparation & Younes \\
12 min 30 & L'explicabilité & Pourquoi et pourquoi pas, garde-fou, conformité & Sohayb \\
14 min & Limites et suite & Ce qui est simulé ; passer à 17\,000 bureaux & Younes \\
\bottomrule
\end{tabularx}
\end{center}

\textbf{Questions attendues.} \emph{Pourquoi pas tout en LLM ?} Mesuré : A contre B.
\emph{L'énergie ?} CodeCarbon par lots, repos soustrait. \emph{Si le classifieur se trompe ?}
Seuil, question, escalade ; courbe à l'appui. \emph{Le LLM invente un taux ?} Il ne reçoit que des
faits, et chaque nombre est contrôlé. \emph{Pourquoi Ministral ?} Le plus petit qui tient le seuil
en descendant, et européen. \emph{Pourquoi ne pas traduire Banking77 ?} L'encodeur est multilingue : on mesure le transfert
au lieu de payer une traduction. \emph{Des mots-clés en 2026 ?} Une part mesurée des demandes pour
presque rien, et ils disent quand ils doutent. \emph{Conformité ?} Pas de notation de solvabilité,
un conseiller décide, transparence, mention légale.

% ───────────────────────── 10. Ce soir ─────────────────────────
\section{Ce qu'on fait ce soir}

Déjà fait : le socle de Sohayb (contrats, \f{reply()}, interface, tests) est sur GitHub ; les
données sont dans \f{data/raw/}, versionnées.

\begin{enumerate}
  \item Lire ce résumé ; valider la répartition, puis ne plus y revenir.
  \item Relire ensemble \f{src/contrats.py} et le squelette, puis les figer.
  \item Rendre le dépôt GitHub privé.
  \item \textbf{Younes} : écrire à Denis Barthou (DeepEirb, modèle vLLM) ; lancer la génération des
        exemples des 8 nouvelles intentions.
  \item \textbf{Yahya} : brancher le parcours carte et ses gabarits sur le squelette.
  \item \textbf{Sohayb} : les 8 clients, le catalogue v0, le rejeu des scénarios.
  \item En filaire : télécharger \f{ministral-3:3b}, \f{qwen3:1.7b}, \f{gemma3:1b} et
        \f{multilingual-e5-small} sur un poste de la salle.
  \item Avant de partir, vérifier J0 : \f{streamlit run app.py} répond.
\end{enumerate}

\end{document}
